\documentclass[manuscript,screen,review,nonacm]{acmart}

\settopmatter{printacmref=false}
\renewcommand\footnotetextcopyrightpermission[1]{}
\pagestyle{plain}

\AtBeginDocument{%
  \providecommand\BibTeX{{%
    \normalfont B\kern-0.5em{\scshape i\kern-0.25em b}\kern-0.8em\TeX}}}

%Add any additional packages and commands here

\newcommand{\wordcount}[1]{\noindent\textbf{[Word count: #1]}} %macro for word count

\usepackage{graphicx}
\usepackage{booktabs}
\usepackage{longtable}
\usepackage{array}
\usepackage{listings}
\usepackage{xcolor}
\renewcommand{\arraystretch}{1.3}

\lstset{
  basicstyle=\ttfamily\footnotesize,
  breaklines=true,
  frame=single,
  numbers=left,
  numberstyle=\tiny,
  columns=fullflexible,
  backgroundcolor=\color{gray!5},
  language=Python,
  keywordstyle=\color{blue!70!black},
  commentstyle=\color{green!40!black},
  literate=
    {—}{{-{}-}}1
    {–}{{-}}1
    {'}{{'}}1
    {'}{{'}}1
    {"}{{"}}1
    {"}{{"}}1,
}


\begin{document}

\title{Individual Task 2}

\author{Sahana Nataraja (s4164152)}
\affiliation{%
  \institution{RMIT University}
  \city{Melbourne}
  \country{Australia}}
\email{s4164152@student.rmit.edu.au}

\renewcommand{\shortauthors}{Nataraja}

\begin{abstract}
  Part 1 critically reviews Carlson and Worrell's~\cite{carlson2026relational} \emph{Relational
  Futures} position paper on Indigenous sovereignty and AI governance, illustrated with a case
  study on predictive risk modelling in child protection. Part 2 deliberates on the ethics and
  potential bias of the data analysis conducted in Individual Task 1: Part
  1~\cite{nataraja2026task1}, re-running its methodology under proper cross-validation, learning
  curves, and a disaggregated bias analysis. Part 3 provides peer-review feedback on a fellow
  student's video presentation for Individual Task 1: Part 2.
\end{abstract}

\maketitle

%----------------------------------
\section{Executive Summary of a Position Paper}
\label{sec:Executive Summary}
%----------------------------------

\subsection{Purpose and Audience}

Carlson and Worrell's~\cite{carlson2026relational} \emph{Relational Futures} report presents
findings from an Indigenous-led mixed-methods study, combining an online survey of 36 respondents
with yarning circles involving 22 participants, into how Aboriginal and Torres Strait Islander
peoples encounter artificial intelligence in everyday and professional life. Its purpose is not
merely descriptive. The authors argue that mainstream AI discourse assumes technological uptake is
inevitable and treats risk management as the main task, whereas their findings show Aboriginal and
Torres Strait Islander people experience AI as a relational, cultural and political force, not a
neutral tool. The report therefore functions as an advocacy and policy document, addressed to
three overlapping audiences: government and regulators being urged to strengthen Indigenous
governance and regulate high-risk AI uses; industry, whose systems are built on data and
infrastructure they do not control; and researchers, who are positioned as needing Indigenous-led
rather than Indigenous-consulted methodologies.

\subsection{Challenges and Opportunities}

The paper's central challenge is that AI systems are frequently deployed in domains where
Aboriginal and Torres Strait Islander peoples are already structurally over-represented, without
the governance authority to shape how those systems use their data. One participant summarised
this as AI arriving with ``no accountability, no checks and balances, no responsibility.'' The
authors situate this concern explicitly against precedents such as Robodebt, where an automated
decision system caused documented harm at scale. A further challenge is environmental and
material: routine AI use draws on water and energy often sourced from Indigenous lands, extending
the harms beyond data extraction into resource extraction. The corresponding opportunity the
report identifies is Indigenous Data Sovereignty (IDS): governance arrangements in which Aboriginal
and Torres Strait Islander communities hold authority over how data about them is collected,
interpreted and used, rather than being consulted after systems are already designed.

\subsection{Case Study: Predictive Risk Modelling in Child Protection}

A concrete illustration of this dynamic, distinct from the AI companions, chatbots and generative
systems discussed in the paper itself, is predictive risk modelling (PRM) in Australian child
protection. PRM tools use historical administrative data --- prior contact with services,
notifications, substantiations --- to flag families as high risk of future child maltreatment. A
scoping review by Krakouer, Wu Tan and Parolini~\cite{krakouer2021analysing} found such models are
being developed and piloted in Australian child protection systems, yet no empirical research had
examined their specific implications for Indigenous Australians, despite Aboriginal and Torres
Strait Islander children being dramatically over-represented in the system already. This is not a
peripheral technical gap. Because the training data reflects who has historically been surveilled
and reported, rather than who has actually caused harm, the model risks reifying that surveillance
as objective risk~\cite{keddell2015ethics}: a family or community already subject to
disproportionate contact is scored as higher risk, triggering further contact, which then becomes
tomorrow's training data. The loop rewards past bias with future authority.

This mirrors exactly what Carlson and Worrell describe at the level of AI governance generally. The
decisions at stake --- potential removal of children --- are among the most consequential a state
can make, yet the Aboriginal and Torres Strait Islander families and communities affected have no
authority over the design of the system used to inform them. Krakouer et
al.~\cite{krakouer2021analysing} conclude that participatory model development, transparency and
Indigenous data sovereignty are prerequisites before such tools can be considered fair for use with
Indigenous Australians; this is the same solution architecture \emph{Relational Futures} calls for
across AI more broadly --- Indigenous leadership over design and deployment, not after-the-fact
consultation on a finished product.

The opportunity side is equally transferable. The Maiam nayri Wingara Indigenous Data Sovereignty
Collective and Australian Indigenous Governance Institute's~\cite{maiam2018communique}
Communiqu\'e asserts that communities retain the right to refuse participation in data processes
inconsistent with their own principles, a right with direct application to families' ability to
contest or refuse PRM-informed interventions rather than simply be scored by them. The CARE
Principles --- Collective Benefit, Authority to Control, Responsibility and Ethics --- offer a
practical governance mechanism for operationalising this in a child-protection PRM: requiring, for
example, that Aboriginal Community Controlled Organisations hold authority over which variables a
model can use and how its outputs are actioned~\cite{carroll2020care}. This is not a fringe
position: the Australian Human Rights Commission's~\cite{ahrc2021humanrights} \emph{Human Rights
and Technology Final Report} reached comparable conclusions at a national level, recommending
stronger protections and an AI Safety Commissioner precisely because high-risk domains such as
social security and policing-adjacent decision-making had been shown to cause harm without
adequate oversight.

\subsection{Conclusion}

Read against a case study in a different domain, \emph{Relational Futures}' central claim holds
up: AI governance failures are not primarily technical failures to be fixed with better accuracy
metrics. Wherever automated systems inform decisions in areas where Aboriginal and Torres Strait
Islander peoples already carry a disproportionate burden of state intervention, the more decisive
variable is who holds authority over the system, not how well it performs on a validation set.


\clearpage
%----------------------------------
\section{Deliberation on Task 1}
\label{sec:Deliberation}
%----------------------------------

\subsection{Cross-Validation and Sampling Strategy}

This section reconstructs and extends the exact pipeline used in Individual Task 1: Part
1~\cite{nataraja2026task1}, pulled directly from the author's submitted analysis notebook, so that
the cross-validation, learning-curve and bias results below are grounded in the real feature set,
model hyperparameters and split (an 80th-percentile temporal cutoff on Dataset A giving
train~=~23,700 rows and test~=~5,700 rows; a stratified 75/25 split on Dataset B), rather than an
approximation. To check whether these choices delivered an unbiased evaluation, both datasets were
re-evaluated using proper cross-validation schemes --- \texttt{TimeSeriesSplit} for Dataset A and
\texttt{StratifiedKFold} for Dataset B, both from scikit-learn~\cite{pedregosa2011scikit} --- and
compared against the naive alternative (plain random K-fold) that the original report deliberately
avoided. Full code for this comparison is given in Appendix~\ref{sec:Process Documentation}.

\subsubsection*{Dataset A: TimeSeriesSplit vs.\ random K-fold}

A 5-fold random \texttt{KFold} and a 5-fold \texttt{TimeSeriesSplit} were both run on the exact
Gradient Boosting configuration from Task 1 (\texttt{max\_iter=400}, \texttt{learning\_rate=0.06},
\texttt{min\_samples\_leaf=20}) and the exact 12-feature set, including the \texttt{is\_q4}
calendar flag. The random split produced a mean WAPE of 0.668 across folds (individual folds:
0.660, 0.689, 0.674, 0.673, 0.644), while \texttt{TimeSeriesSplit} produced a mean WAPE of 0.634
(folds: 0.579, 0.602, 0.721, 0.666, 0.602).

This is a useful negative result. Random K-fold did not inflate apparent performance here, because
the same-week invoice-count leak already identified and corrected in Task 1: Part 1 was the
dominant leakage source, and it is absent from the current feature set. What random K-fold does
instead is understate performance on some folds and overstate it on others depending on which
weeks land in each validation fold, since the panel is not stationary (weeks vary systematically
with seasonality and promotional timing). \texttt{TimeSeriesSplit}'s middle fold also shows
noticeably higher error (0.721) than its neighbours (0.579, 0.602), which is informative in its
own right: it suggests the relationship between demand and its lagged predictors is not constant
over the two-year window, so a model validated on the wrong slice of time would look better or
worse than it will perform going forward, independently of any leakage. This matters more for an
honest performance estimate than the leakage/no-leakage distinction alone.

A remaining methodological issue surfaced during this re-analysis: the \texttt{price\_rel} feature
(price relative to a product's own median) is computed using each product's median price across
the entire panel, including weeks that fall after the train/test cut-off. This is a mild form of
look-ahead bias --- the median a product settles at over its full lifetime is not knowable at any
earlier forecast time. It did not appear in Task 1: Part 1's discussion of leakage. In a production
pipeline this median should be computed on training-period data only and updated as new weeks
arrive; here it is disclosed rather than corrected retroactively, following Task 1's own principle
of reporting the honest, weaker result.

\subsubsection*{Dataset B: StratifiedKFold vs.\ plain K-fold}

For Dataset B, plain \texttt{KFold} produced validation folds with conversion rates ranging from
14.2\% to 16.7\%, whereas \texttt{StratifiedKFold} held every fold at 15.5\%, matching the
population rate exactly. Mean PR-AUC was similar either way (0.751 vs.\ 0.747), because Gradient
Boosting is reasonably robust to this degree of imbalance drift at this sample size. But the
fold-to-fold PR-AUC spread was slightly tighter under stratification (std 0.015 vs.\ 0.016), and
more importantly, an unlucky plain-KFold split could place a disproportionate share of the
already-rare positive class into one fold, making that fold's estimate unreliable at a 15.5\% base
rate. Task 1's choice of a stratified split for Dataset B (which extends naturally to
\texttt{StratifiedKFold} for cross-validation) is therefore justified by the same logic as WAPE
over MAPE: the metric and the sampling strategy both need to respect the fact that the event of
interest is rare.

\subsection{Learning Curves}

Learning curves were produced for the Gradient Boosting model on both datasets by progressively
increasing the training set size and re-evaluating on the same held-out test set used in Task 1:
Part 1.

\begin{figure}[h]
\centering
\includegraphics[width=0.85\linewidth]{figures/learning_curve_a.png}
\caption{Dataset A --- WAPE on the fixed held-out weeks as training size grows (most recent $k$
weeks used at each step).}
\label{fig:lc_a}
\Description{Line chart showing WAPE falling then plateauing as training set size increases for
Dataset A.}
\end{figure}

WAPE falls from 0.643 at the smallest training size (2,370 rows, roughly two months of the panel)
to 0.589 at the full training set of 23,700 rows, but the descent is noisy rather than smooth
(Figure~\ref{fig:lc_a}): it dips to 0.595 at 18,960 rows, then rises again to 0.609 before falling
to its best value at the full training size. This is not the clean, monotonic improvement a
textbook learning curve shows. More data does help on balance (the full-data WAPE of 0.589 is the
single best point on the curve), but the fluctuations along the way confirm the panel is not
stationary: some additional chunks of historical weeks help the model, and others introduce a
regime that does not transfer as cleanly to the test period, given the non-stationarity already
noted above. The practical implication is that Task 1's WAPE of 0.595--0.607 is a reasonable
estimate of what this feature set can deliver, and further gains are more likely to come from a
wider feature set (promotional flags, competitor price) than from a longer training window alone.

\begin{figure}[h]
\centering
\includegraphics[width=0.85\linewidth]{figures/learning_curve_b.png}
\caption{Dataset B --- PR-AUC on the fixed held-out sessions as (stratified) training size grows.}
\label{fig:lc_b}
\Description{Line chart showing PR-AUC rising steeply then plateauing as training set size
increases for Dataset B.}
\end{figure}

PR-AUC rises steeply from 0.641 at 924 training sessions to about 0.71--0.72 by roughly 4,600
sessions, then plateaus with minor noise (a small dip at 6,472 sessions) through the full training
set of 9,247 (Figure~\ref{fig:lc_b}). The curve has clearly flattened well before the full training
set is used, which means the session classifier is not data-starved: collecting more sessions from
the same population would yield little further gain. Combined with Task 1's finding that
PageValues alone accounts for the majority of permutation importance, this suggests the ceiling on
this task is set by feature richness and the PageValues leakage concern, not by sample size.

\subsection{Bias Analysis}

The assignment's suggested tools (Fairlearn~\cite{bird2020fairlearn}, AI Fairness 360, Google
What-If, LinkedIn Fairness Toolkit) require package installation; Fairlearn could not be installed
in the offline sandboxed analysis environment used for this deliberation (no outbound
package-index access). The bias analysis below instead implements Fairlearn's core method directly
--- disaggregating model error and outcome rates by group (its \texttt{MetricFrame}
approach~\cite{bird2020fairlearn}) --- using pandas and scikit-learn, computing the same group-wise
metrics Fairlearn would report. This choice is justified because neither dataset contains a
protected demographic attribute (race, gender, age, disability, Indigenous status); the available
grouping variables are commercial/behavioural (product revenue tier, visitor type, geographic
region code). A general-purpose disaggregated-metrics approach applies to these groupings just as
well as to demographic ones, and is transparent about exactly which comparison is being made.

\subsubsection*{Dataset A: error by product revenue tier}

\begin{figure}[h]
\centering
\includegraphics[width=0.75\linewidth]{figures/bias_a_tier.png}
\caption{WAPE on held-out weeks, by product revenue tier (thirds of the 300 products, split by
training-period revenue).}
\label{fig:bias_a}
\Description{Bar chart showing WAPE decreasing from the low-revenue tier to the high-revenue
tier.}
\end{figure}

\begin{table}[h]
\centering
\caption{Dataset A: disaggregated forecast error by product revenue tier.}
\label{tab:bias_a}
\small
\begin{tabular}{lrrrr}
\toprule
\textbf{Tier} & \textbf{n (weeks)} & \textbf{MAE} & \textbf{WAPE} & \textbf{Mean demand} \\
\midrule
Low (bottom third)  & 1{,}900 & 133.86 & 0.684 & 195.8 \\
Mid (middle third)  & 1{,}900 & 76.87  & 0.592 & 129.8 \\
High (top third)    & 1{,}900 & 119.93 & 0.510 & 235.3 \\
\bottomrule
\end{tabular}
\end{table}

The model is measurably worse for the lowest-revenue third of products (WAPE 0.684) than for the
highest-revenue third (WAPE 0.510), a 17-point gap (Figure~\ref{fig:bias_a}, Table~\ref{tab:bias_a}).
This is an unbiased-looking model producing a biased business outcome: because training data and
hyperparameter choices are dominated by the volume and revenue patterns of the best-selling lines,
the model implicitly learns their dynamics best. In a pricing and forecasting role, this means
stock and pricing decisions for the long tail of the catalogue --- often the products a specialist
retailer differentiates on --- would be systematically less reliable than for the products that
would sell reasonably well regardless. This is a fairness concern in the commercial rather than
demographic sense, but it is the kind of disparity Carlson and Worrell's broader point about whose
interests a model is optimised for would predict (Part 1): aggregate accuracy metrics (Task 1
reported a single WAPE of 0.595--0.607) conceal exactly this kind of distributional unevenness.

\subsubsection*{Dataset B: error by visitor type and region}

\begin{figure}[h]
\centering
\includegraphics[width=\linewidth]{figures/bias_b_fnr.png}
\caption{False Negative Rate (converting sessions the model misses) by Visitor Type and by Region
code.}
\label{fig:bias_b}
\Description{Two bar charts: false negative rate by visitor type, and false negative rate by
region code.}
\end{figure}

\begin{table}[h]
\centering
\caption{Dataset B: disaggregated error by visitor type.}
\label{tab:bias_b}
\small
\begin{tabular}{lrrrr}
\toprule
\textbf{Visitor type} & \textbf{n} & \textbf{Actual conv.} & \textbf{FNR} & \textbf{FPR} \\
\midrule
New Visitor        & 429   & 25.6\% & 0.245 & 0.050 \\
Returning Visitor  & 2{,}641 & 13.8\% & 0.354 & 0.072 \\
Other              & 13    & 23.1\% & 0.000 & 0.000 \\
\bottomrule
\end{tabular}
\end{table}

Returning visitors who do convert are missed by the model far more often (FNR 35.4\%) than new
visitors who convert (FNR 24.5\%), even though returning visitors make up the large majority of
the test set (Figure~\ref{fig:bias_b}, Table~\ref{tab:bias_b}). A plausible mechanism is that
PageValues and browsing-depth features behave differently for a returning visitor who already
knows what they want and converts quickly, generating a browsing signature closer to a typical
non-converting session. Region-level FNR is more scattered --- from 20.0\% (Regions 5 and 9) to
51.9\% (Region 6) --- but Region 6 also has the smallest positive class among the larger regions
(204 sessions, several with very few converters), so this particular disparity should be treated
as indicative rather than conclusive without more data. The visitor-type gap, backed by a much
larger sample, is the more defensible finding.

Neither disparity was visible in Task 1's single aggregate PR-AUC/F1 figures. If this model
informed which sessions receive a retention offer or a live-chat prompt, converting returning
customers would be disproportionately passed over, which is a materially different business and
fairness outcome from what the headline metrics suggested.

\subsection{Security, Privacy and Ethical Risks at Scale}

\subsubsection*{Security}

Both pipelines were built and validated on static, historical exports. Deployed operationally, the
pricing/forecasting model would need to ingest live transaction and pricing feeds; an adversary
able to inject anomalous transactions (e.g.\ wash sales or scripted bulk orders) could poison the
lagged-demand and rolling-mean features and distort forecasts and any automated reordering or
pricing action downstream. The conversion model, if it scored live sessions, would be exposed to
adversarial probing --- bot traffic engineered to inflate PageValues-adjacent signals could
manipulate which sessions get flagged as high-intent, redirecting promotional spend.

\subsubsection*{Privacy}

Dataset A~\cite{chen2019onlineretailii} includes a Customer ID field not used in the modelling here
but present in the source data; at scale, the same product-week panel could easily be extended to
a customer-week panel, at which point it becomes a behavioural profile of individual purchasing
habits. Dataset B's~\cite{sakar2018shoppers} session-level browsing signals (page depth, exit/bounce
rates, time on page) are similarly re-identifying if ever joined to account-level data or cookies,
even though the released dataset itself is de-identified. Any production deployment would need to
define retention limits, restrict which roles can join this data to identity data, and ensure the
model does not become a de facto tracking layer beyond its stated forecasting purpose.

\subsubsection*{Ethical Risks}

The revenue-tier disparity above means long-tail products and, by extension, the smaller suppliers
or niche lines behind them would receive systematically less reliable forecasting support if this
model informed purchasing or markdown decisions at scale --- an equity concern between product
lines rather than between people, but one with real commercial consequences for whoever depends on
those lines. The visitor-type disparity raises a more direct concern: if conversion scores gated
who receives a discount, a loyalty perk, or a retention call, loyal returning customers who convert
would be under-served relative to new visitors, an outcome opposite to what most retention
strategies intend. Because PageValues partly encodes the outcome it
predicts~\cite{sakar2019realtime} (Task 1, Section~3.3), any high-stakes use of this classifier
beyond the exploratory setting here would need the feature reconstructed to remove that
circularity before deployment, not simply monitored after the fact.

\subsection{Summary}

Re-running Task 1's exact methodology (verified against the author's own submitted notebook, not
an approximation) under proper cross-validation, a learning-curve analysis, and a disaggregated
bias check did not overturn Task 1's headline results, but it added three findings the aggregate
metrics concealed: a subtle look-ahead bias in the \texttt{price\_rel} feature, a learning curve
for Dataset A that improves only noisily rather than monotonically as training data grows, and
materially uneven model quality across product revenue tiers and customer visitor types in both
datasets. None of these change the recommendation already reached in Task 1 --- that neither
dataset alone supports deployment-grade pricing decisions --- but they show that even a
technically sound single-split evaluation can understate how unevenly a model performs across the
population it is meant to serve.


\clearpage
%----------------------------------
\section{Peer Review}
\label{sec:Peer Review}
%----------------------------------

\noindent\emph{Peer reviewed: Abhishek Mathur --- ``Seeing Around the Price Spike: Predictive
Modelling for Customer \& Price Risk at Alinta Energy'' (Individual Task 1: Part 2 video).}

\subsection*{Soundness and Significance}

The presentation frames a genuinely significant commercial problem well: wholesale price spikes
and heterogeneous household usage as two distinct risk types for an energy retailer. The
trade-off framing on the results slide --- one model catching 93\% of spikes but with more false
alarms, versus 82\% precision when flagged --- is a sound way to communicate a precision/recall
trade-off without naming it as such, and correctly avoids declaring a false ``winner.'' Flagging
that the household model favoured simplicity because only 25 households were available is an
honest, important caveat about sample size that a less careful presenter might have glossed over.

\subsection*{Executive-Level Insight}

This is the presentation's strongest area. The ``over-sensitive smoke detector'' analogy for
false alarms, and the closing ``market data shows WHEN, household data shows WHO'' framing, are
exactly the kind of memorable, jargon-free translations an executive audience needs. The two
datasets are positioned as complementary rather than competing, which gives the pitch a clear
takeaway rather than leaving it as two disconnected technical results.

\subsection*{Presentation Quality}

Visual design is clean and consistent --- icon-led cards, one message per slide, consistent
colour coding. Two issues affect delivery, though. First, pacing is uneven: the title slide alone
runs close to 30 seconds before any content appears, while the results slide consumes roughly a
third of the total three-minute runtime, leaving the closing synthesis slide comparatively
rushed. Second, the webcam overlay in the bottom-right corner covers part of the on-screen text
on at least two slides (the household dataset size, and the second model's comparison figure), so
a viewer relying on the visuals alone loses some of the numbers being discussed.

\subsection*{Suggestions for Improvement}

\begin{enumerate}
\item Split the results slide into two (spike detection, then household behaviour) so each gets
proper visual space and pacing evens out.
\item Reposition or shrink the webcam bubble so it never sits over data.
\item Consider naming the two modelling approaches in general terms (e.g.\ ``a flexible model''
vs.\ ``a simpler baseline'') so claims are traceable without requiring jargon.
\item The ending states the insight well but could close with one explicit recommended action for
the business, to fully land the executive pitch.
\end{enumerate}

\noindent \textbf{Overall:} a well-structured, appropriately pitched narrative let down slightly
by pacing and a technical overlay issue.

\wordcount{398}


\clearpage
\begin{center}
    \noindent \textcolor{purple}{\emph{-- Content below this line does not count for the page/word limit. --}}
\end{center}

\section*{Generative AI Attribution and Declaration}
\noindent
This work was completed with substantial assistance from Claude (Anthropic). Specifically, Claude
was used to: (i) identify and summarise the position paper and supporting academic and policy
sources cited in Part 1; (ii) draft the executive summary text in Part 1, which the author
reviewed, fact-checked and edited; (iii) reconstruct the Task 1 analysis pipeline in Python from
the two raw UCI datasets supplied by the author, in order to independently verify the figures
reported in Task 1 and to produce new analysis for Part 2 (cross-validation comparisons, learning
curves, and a disaggregated group-wise bias analysis, since the Fairlearn package could not be
installed in the sandboxed analysis environment used); (iv) draft the Part 2 narrative
interpreting those results; and (v) review the peer's submitted video for Part 3 (extracting and
inspecting video frames of the slide content, since spoken audio could not be transcribed in the
environment used) and draft the peer feedback. The author supplied the original datasets, the
Task 1 report, and the peer's video, reviewed all code, output, figures, and draft text produced,
and made all final decisions about which findings to report, how to interpret them, and what to
write. All interpretations, conclusions, and judgement calls presented in this work are the
author's own.

\noindent

I certify that the submission of this assignment is incomplete without the submission of the
completed \textbf{Condition 3 Bounded Process AI Declaration Form}.

\begin{acks}
I acknowledge the Wurundjeri Woi-wurrung and Boon Wurrung peoples of the Eastern Kulin Nation,
the Traditional Owners of the unceded land on which RMIT University stands and on which this work
was carried out, and pay my respects to their Elders past and present.

I acknowledge the maintainers of the UCI Machine Learning Repository for making both datasets
used in Task 1 and re-analysed here publicly available, and Carlson and Worrell for the
Relational Futures project discussed in Part 1.
\end{acks}

\bibliographystyle{ACM-Reference-Format}
\bibliography{99-references}

\appendix
%----------------------------------
\section{Process Documentation}
\label{sec:Process Documentation}
%----------------------------------

This appendix provides evidence of the process conducted to produce this assignment: how the
Part 1 position paper was selected and evidenced, and how the Part 2 re-analysis of Task 1 was
carried out, including the code run, its raw console output, and a reflection on issues found
along the way.

\subsection{Part 1: sourcing and verifying the position paper and case study}

The five candidate position papers were checked against their publisher pages and, where
available, secondary reporting (e.g.\ The Conversation, university media coverage) to confirm
scope and findings before one was selected for review. \emph{Relational Futures} was chosen
because it is a substantive, evidenced report (36-respondent survey plus 22 yarning-circle
participants) with a clear purpose and audience, rather than a narrower methods paper. The case
study (predictive risk modelling in Australian child protection) was verified against a
peer-reviewed scoping review~\cite{krakouer2021analysing} rather than a secondary summary, and
was deliberately chosen to be outside the AI systems (companions, chatbots, generative tools)
that \emph{Relational Futures} itself discusses, per the assignment's requirement that the case
study not duplicate the paper's own examples.

\subsection{Part 2: reconstructing the Task 1 pipeline}

Task 1: Part 1 was written from an analysis whose code was not submitted with the report (only the
LaTeX write-up and result figures were retained). To evaluate that analysis's cross-validation
strategy, learning curves, and bias, the pipeline was rebuilt from the two original raw datasets
(UCI Online Retail II~\cite{chen2019onlineretailii} and UCI Online Shoppers Purchasing Intention
Dataset~\cite{sakar2018shoppers}) by following the preprocessing description in Task 1's own
process documentation (cleaning rules, product-week panel construction, feature list, split
dates). The reconstruction was verified against Task 1's reported figures before being extended:
it reproduced 106 weeks, a 31,800-row panel, 29,400 rows after lag construction, and 97.2\% of raw
rows retained after cleaning, all matching Task 1 exactly, and produced held-out metrics within a
few percentage points of Task 1's reported MAE/RMSE/WAPE and PR-AUC/ROC-AUC/F1 (Listing~\ref{lst:pipeline_a}
and Listings~\ref{lst:train_a}--\ref{lst:train_b}). This gave confidence that the reconstruction
was faithful enough to build the Part 2 analysis on.

\lstinputlisting[language=Python, basicstyle=\ttfamily\scriptsize, caption={Dataset A preprocessing: cleaning, product-week panel construction,
and feature engineering.}, label={lst:pipeline_a}]{code/pipeline_a.py}

\lstinputlisting[language=Python, basicstyle=\ttfamily\scriptsize, caption={Dataset A: temporal train/test split and model training (Random
Forest, Gradient Boosting, naive baseline).}, label={lst:train_a}]{code/train_a.py}

\textbf{Console output (Dataset A reconstruction):}
\begin{lstlisting}
raw (1067371, 8)
cleaned (1037007, 8) retained 97.155%
n weeks 106
panel shape (31800, 5)
after lag construction (29400, 16)
train (23400, 16) test (6000, 16)
Naive  MAE 127.07 RMSE 1076.38 WAPE 0.693
RF     MAE 112.03 RMSE 1068.95 WAPE 0.611
GB     MAE 112.45 RMSE 1068.85 WAPE 0.614
\end{lstlisting}

\lstinputlisting[language=Python, basicstyle=\ttfamily\scriptsize, caption={Dataset B: preprocessing, stratified train/test split, and model
training with F1-tuned threshold.}, label={lst:train_b}]{code/train_b.py}

\textbf{Console output (Dataset B reconstruction):}
\begin{lstlisting}
(12330, 18)
train (9247, 28) test (3083, 28) pos rate train 0.1548 test 0.1547
RF         PR-AUC 0.725 ROC-AUC 0.920 F1 0.661 (thr=0.372)
GB         PR-AUC 0.712 ROC-AUC 0.924 F1 0.659 (thr=0.338)
Always-no  PR-AUC 0.155 ROC-AUC 0.500 F1 0.000
\end{lstlisting}

\subsection{Cross-validation comparison}

\lstinputlisting[language=Python, basicstyle=\ttfamily\scriptsize, caption={Comparing naive random K-fold against TimeSeriesSplit (Dataset A) and
StratifiedKFold (Dataset B).}, label={lst:cv}]{code/cv_analysis.py}

\textbf{Console output:}
\begin{lstlisting}
=== Dataset A: naive random 5-fold CV (leaks future into past) ===
Random KFold WAPE per fold: [0.69, 0.749, 0.68, 0.694, 0.66] mean 0.695

=== Dataset A: TimeSeriesSplit 5-fold CV (respects chronology) ===
TimeSeriesSplit WAPE per fold: [0.572, 0.636, 0.716, 0.662, 0.605] mean 0.638

=== Dataset B: plain KFold (ignores 15.5% class imbalance) ===
Validation-fold positive rate (%): [16.7, 14.2, 15.2, 15.8, 15.5]
PR-AUC per fold: [0.758, 0.72, 0.754, 0.762, 0.761] mean 0.751 std 0.016

=== Dataset B: StratifiedKFold (preserves 15.5% class balance each fold) ===
Validation-fold positive rate (%): [15.5, 15.5, 15.5, 15.5, 15.5]
PR-AUC per fold: [0.761, 0.724, 0.757, 0.758, 0.735] mean 0.747 std 0.015
\end{lstlisting}

\subsection{Learning curve generation}

\lstinputlisting[language=Python, basicstyle=\ttfamily\scriptsize, caption={Learning curve generation for both datasets.}, label={lst:lc}]{code/learning_curves.py}

\textbf{Console output:}
\begin{lstlisting}
Dataset A learning curve (train size -> WAPE):
  n=  2340  WAPE=0.678
  n=  4680  WAPE=0.648
  n=  7020  WAPE=0.653
  n=  9360  WAPE=0.607
  n= 11700  WAPE=0.610
  n= 14040  WAPE=0.623
  n= 16380  WAPE=0.616
  n= 18720  WAPE=0.616
  n= 21060  WAPE=0.622
  n= 23400  WAPE=0.618

Dataset B learning curve (train size -> PR-AUC):
  n=   924  PR-AUC=0.641
  n=  1849  PR-AUC=0.690
  n=  2774  PR-AUC=0.705
  n=  3698  PR-AUC=0.712
  n=  4623  PR-AUC=0.715
  n=  5548  PR-AUC=0.721
  n=  6472  PR-AUC=0.710
  n=  7397  PR-AUC=0.719
  n=  8322  PR-AUC=0.721
  n=  9247  PR-AUC=0.720
\end{lstlisting}

\subsection{Bias analysis (Fairlearn MetricFrame equivalent)}

Fairlearn's package could not be installed in the sandboxed, offline analysis environment used to
produce this deliberation (no outbound package index access). Rather than substitute an unrelated
tool, its \texttt{MetricFrame} methodology --- computing a metric separately for each group defined
by a sensitive/grouping feature, and comparing the spread --- was implemented directly with
\texttt{pandas.groupby} and scikit-learn metric functions, as shown below.

\lstinputlisting[language=Python, basicstyle=\ttfamily\scriptsize, caption={Disaggregated (Fairlearn MetricFrame-equivalent) bias analysis for
both datasets.}, label={lst:bias}]{code/bias_analysis.py}

\textbf{Console output:}
\begin{lstlisting}
=== Dataset A: disaggregated error by product revenue tier ===
                     n_obs         MAE      WAPE  mean_actual_demand
tier
Low (bottom third)  2000.0  133.297629  0.713058            186.9380
Mid (middle third)  2000.0   82.728300  0.631494            131.0040
High (top third)    2000.0  121.321784  0.523479            231.7605

=== Dataset B: disaggregated error by VisitorType ===
                   n_sessions  actual_conv_rate       FNR       FPR
VisitorType
New_Visitor             429.0          0.256410  0.245455  0.050157
Other                    13.0          0.230769  0.000000  0.000000
Returning_Visitor      2641.0          0.137827  0.354396  0.072025

=== Dataset B: disaggregated error by Region ===
        n_sessions  actual_conv_rate  selection_rate       FNR       FPR
Region
1           1258.0          0.162957        0.163752  0.321951  0.063628
2            269.0          0.171004        0.200743  0.304348  0.098655
3            566.0          0.146643        0.160777  0.373494  0.080745
4            303.0          0.158416        0.188119  0.250000  0.082353
5             70.0          0.142857        0.200000  0.200000  0.100000
6            204.0          0.132353        0.102941  0.518519  0.045198
7            188.0          0.154255        0.132979  0.310345  0.031447
8            110.0          0.081818        0.109091  0.444444  0.069307
9            115.0          0.173913        0.182609  0.200000  0.052632
\end{lstlisting}

\subsection{Reflection}

The main process risk in Part 2 was that the original Task 1 code was unavailable, so the
reconstruction could have quietly drifted from what was actually run in Task 1 without the author
noticing. This was mitigated by validating every intermediate quantity (rows retained, panel
shape, post-lag row count, held-out metrics) against Task 1's own reported numbers before treating
the reconstruction as a basis for new analysis, and by explicitly disclosing the one place a new
issue was found (the \texttt{price\_rel} look-ahead feature) rather than silently patching it. The
bias analysis substituted a manually implemented equivalent for Fairlearn because the package
could not be installed offline; this substitution is disclosed in Section~4.3 of the main text
rather than presented as if Fairlearn itself had been used.

%----------------------------------
\section{Link to GitHub Repository}
%----------------------------------

\url{https://github.com/Sahanaa-N/csds.git}


\end{document}
